% TEMPLATE-MILC-v3.tex - plantilla maestra UNICA de las guias MILC v3.
%
% La usan las tres familias de guias, cada una con su driver:
%   · Grados 6-11  -> scripts/build-guias-g11.py
%   · Bebras       -> scripts/build-guias-bebras.py
%   · Semillero    -> scripts/build-guias-semillero.py
%
% Antes habia tres copias casi identicas (~400 lineas iguales, 6 distintas):
% cada mejora habia que aplicarla tres veces, y en la practica no se hacia
% (el motor de recursos vivia solo en dos de ellas). Lo que varia entre
% familias esta parametrizado en 6 placeholders que cada driver llena
% (se nombran aqui SIN su sintaxis real: el builder reemplaza por texto plano,
% tambien dentro de los comentarios, y un valor multilinea rompe el preambulo):
%
%   HEADER_IZQ        encabezado izquierdo de pagina
%   HEADER_DER        encabezado derecho de pagina
%   PORTADA_ETIQUETA  rotulo sobre el numero grande (GRADO/MOMENTO/MODULO)
%   PORTADA_SERIE     linea de serie de la portada
%   PORTADA_ID        identificador de la guia en la portada
%   PORTADA_CIRCULO   el circulito de grado (°), o vacio si no aplica
%
% El resto de claves las documenta content/guias/_SCHEMA.md. El script valida
% que no queden placeholders sin reemplazar antes de escribir el archivo final.

\documentclass[11pt,letterpaper]{article}
\usepackage[margin=2.0cm]{geometry}
\usepackage[table]{xcolor}
\usepackage{fontspec}
\usepackage[spanish]{babel}
\usepackage{tikz}
% Librerías TikZ para los diagramas de las guías (bloque `recursos:`):
%  · babel  → babel en español vuelve activos caracteres como «>», y TikZ los usa
%             en sus opciones (>=stealth, ->). Sin esta librería, cualquier
%             diagrama con flechas revienta con «\language@active@arg> has an
%             extra }». Debe ir ANTES de que se dibuje nada.
%  · positioning        → sintaxis «below left=2pt and 3pt».
%  · decorations.pathreplacing → llaves ({brace}) para señalar rangos.
\usetikzlibrary{babel,positioning,decorations.pathreplacing}
\usepackage{graphicx}
% Circuitos: el trabajo de electrónica se hace hoy con micro:bit y sus sensores
% integrados (MakeCode, sin cableado), así que no se necesitan esquemáticos.
% Si algún día entran montajes con protoboard, basta descomentar esta línea:
% los diagramas .tex/.tikz declarados en `recursos:` ya se incrustan solos.
% \usepackage{circuitikz}
\usepackage[most]{tcolorbox}
\usepackage{tabularx}
\usepackage{array}
\usepackage{fancyhdr}
\usepackage{titlesec}
\usepackage[document]{ragged2e}  % bandera a la derecha en todo el cuerpo
\usepackage{enumitem}
\usepackage{microtype}

% Helvetica Neue no trae la flecha U+2192, asi que cada «→» escrito literalmente
% en el YAML se compilaba a NADA, en silencio (462 flechas en 61 guias: rutas del
% tipo «paleta → tipografia → lockup» salian rotas). Mapeamos el caracter a su
% version matematica, que si tiene glifo. Asi el autor puede seguir escribiendo
% «→» comodamente en el YAML.
\usepackage{newunicodechar}
\newunicodechar{→}{\ensuremath{\rightarrow}}
% Mismo problema con los simbolos de marca y vinetas: el «✓» de «una palomita ✓»
% salia como cuadrito vacio en el PDF de todas las guias que lo usaban.
\usepackage{pifont}
\newunicodechar{✓}{\ding{51}}
\newunicodechar{✗}{\ding{55}}
\newunicodechar{✔}{\ding{52}}
\newunicodechar{•}{\textbullet}
\newunicodechar{≥}{\ensuremath{\geq}}
\newunicodechar{≤}{\ensuremath{\leq}}

\setmainfont{Helvetica Neue}
\setsansfont{Helvetica Neue}
\setlength{\parindent}{0pt}
\setlength{\parskip}{6pt}
% Sin particion de palabras: en 6.º-7.º una palabra cortada por guion al final
% de linea («Comuni-cacion») frena la lectura mucho mas que una linea corta.
\hyphenpenalty=10000
\exhyphenpenalty=10000
\pretolerance=2000
\tolerance=3000
\setlength{\emergencystretch}{3em}
\linespread{1.10}
\renewcommand{\arraystretch}{1.25}
\setlist{leftmargin=5mm,itemsep=1mm,topsep=1mm}
\newcolumntype{Y}{>{\RaggedRight\arraybackslash}X}

\definecolor{milcMagenta}{HTML}{E6007E}
\definecolor{milcVino}{HTML}{5A0038}
\definecolor{milcTurquesa}{HTML}{0093A5}
\definecolor{milcVerde}{HTML}{58A923}
\definecolor{milcMostaza}{HTML}{E5B400}
\definecolor{milcOcre}{HTML}{B86600}
\definecolor{milcNegro}{HTML}{241816}
\definecolor{milcGris}{HTML}{F7F7F4}
\definecolor{milcLinea}{HTML}{E8E2DD}
\definecolor{milcGradePrimary}{HTML}{0066FF}
\definecolor{milcGradeDark}{HTML}{003D99}
\definecolor{milcGradeSoft}{HTML}{D6E8FF}
\definecolor{milcGradeAccent}{HTML}{CFE7FF}
\definecolor{milcGradeText}{HTML}{FFFFFF}
\color{milcNegro}

\pagestyle{fancy}
\fancyhf{}
\lhead{\small Tecnología e Informática · Grado 11}
\rhead{\small Guía 29 · MILC v3}
\cfoot{\small\thepage}
\renewcommand{\headrulewidth}{0pt}

\newtcolorbox{titlebox}[1]{
  enhanced,colback=#1,colframe=#1,arc=7mm,boxrule=0pt,
  left=7mm,right=7mm,top=3mm,bottom=3mm,
  before skip=8pt,after skip=8pt,
  fontupper=\sffamily\bfseries\Large\color{white}
}

\newtcolorbox{softbox}[3]{
  enhanced,colback=#2,colframe=#1,borderline west={4pt}{0pt}{#1},
  arc=2mm,boxrule=.4pt,left=4mm,right=4mm,top=3mm,bottom=3mm,
  before skip=6pt,after skip=8pt,title={#3},coltitle=white,
  colbacktitle=#1,fonttitle=\sffamily\bfseries,fontupper=\RaggedRight,
  attach boxed title to top left={xshift=3mm,yshift=-2mm},
  boxed title style={arc=2mm, boxrule=0pt}
}

\newtcolorbox{drawbox}[2]{
  enhanced,colback=white,colframe=#1,arc=4mm,boxrule=1pt,
  left=5mm,right=5mm,top=4mm,bottom=4mm,before skip=8pt,after skip=8pt,
  title={#2},coltitle=white,colbacktitle=#1,fonttitle=\sffamily\bfseries,
  fontupper=\RaggedRight,
  attach boxed title to top left={xshift=4mm,yshift=-2mm},
  boxed title style={arc=3mm, boxrule=0pt}
}

\newtcolorbox{infoband}[1]{
  enhanced,colback=#1!8,colframe=#1!50,arc=2mm,boxrule=.5pt,
  left=4mm,right=4mm,top=2mm,bottom=2mm,before skip=4pt,after skip=4pt,
  fontupper=\small\RaggedRight
}

% Puente narrativo entre secciones (contrato editorial Conéctate).
% Da continuidad: "Acabas de X. Ahora vas a Y porque Z."
% Visualmente: banda izquierda en color del grado, texto en itálica pequeña.
\newtcolorbox{puentebox}{
  enhanced,colback=milcGradeSoft,colframe=milcGradeDark,
  borderline west={3pt}{0pt}{milcGradeDark},
  arc=0mm,boxrule=0pt,left=5mm,right=4mm,top=2.5mm,bottom=2.5mm,
  before skip=4pt,after skip=8pt,
  fontupper=\itshape\small\color{milcNegro}\RaggedRight
}
% La flecha va en modo matematico a proposito: Helvetica Neue NO tiene el glifo
% U+2192, asi que \textrightarrow se compilaba a NADA («Missing character: There
% is no → in font Helvetica Neue») y el puente salia sin flecha. En modo
% matematico la flecha viene de la fuente matematica y si se dibuja.
\newcommand{\puente}[1]{%
  \begin{puentebox}%
    \textcolor{milcGradeDark}{$\rightarrow$}\hspace{2mm}#1%
  \end{puentebox}%
}

\newcommand{\linead}[1]{\noindent\color{milcLinea}\rule{#1}{0.5pt}\color{milcNegro}}
\newcommand{\respuesta}[1]{\par\vspace{2mm}\foreach \n in {1,...,#1}{\noindent\color{milcLinea}\rule{\linewidth}{0.5pt}\par\vspace{5mm}}\color{milcNegro}}
\newcommand{\checkbox}{\textcolor{milcGradeDark}{\fbox{\phantom{x}}}\hspace{5pt}}

\newcommand{\phasecard}[4]{
\begin{tcolorbox}[enhanced,colback=#3,colframe=#2,arc=3mm,boxrule=.5pt,height=3.05cm,valign=center,halign=center,left=1.5mm,right=1.5mm,top=2mm,bottom=2mm]
{\sffamily\bfseries\fontsize{22}{24}\selectfont\textcolor{#2}{#1}}\par
{\sffamily\bfseries\small #4}
\end{tcolorbox}
}

% ── Piezas del rediseno 2026 (legibilidad 6.º-11.º) ───────────────────
% Banda de estacion: reemplaza el titlebox macizo. Titulo a la izquierda,
% dato practico (tiempo/modalidad) a la derecha, en una sola linea.
\newcommand{\milcestacion}[2]{%
\begin{tcolorbox}[enhanced,colback=#1,colframe=#1,arc=2mm,boxrule=0pt,
  left=5mm,right=5mm,top=2.6mm,bottom=2.6mm,before skip=11pt,after skip=7pt]
{\sffamily\bfseries\fontsize{15}{18}\selectfont\color{white}#2}
\end{tcolorbox}}

% Tarjeta de paso numerada: sustituye las tablas «Pilar / Aplicacion», que
% eran formato de consulta docente, no de estudiante. El numero va en un
% circulo sobre el borde para que la secuencia se lea de un vistazo.
\newcommand{\milcpaso}[3]{%
\begin{tcolorbox}[enhanced,colback=white,colframe=milcLinea,arc=1.5mm,boxrule=.9pt,
  left=14mm,right=4mm,top=3mm,bottom=3mm,before skip=4pt,after skip=4pt,
  fontupper=\RaggedRight,
  overlay={\node[anchor=north west,circle,fill=#2,text=white,inner sep=0pt,
    minimum size=7.4mm,font=\sffamily\bfseries\small]
    at ([xshift=3.6mm,yshift=-3.2mm]frame.north west) {#1};}]
#3
\end{tcolorbox}}

% Casilla marcable de tamano util con lapiz (la anterior era \fbox{\phantom{x}},
% de alto variable segun el texto que la seguia).
\newcommand{\milcmarca}{\framebox[3.8mm]{\rule{0pt}{3.4mm}}}

% Fila marcable de lista suelta (checks de la estacion 1).
\newcommand{\milccheckrow}[1]{%
\par\vspace{1.8mm}\noindent
\makebox[6.4mm][l]{\raisebox{-.55mm}{\textcolor{milcNegro}{\milcmarca}}}%
\parbox[t]{\dimexpr\linewidth-6.4mm\relax}{\RaggedRight #1}\par}

% Celda marcable dentro de la rubrica (tabla de autoevaluacion). Misma
% sangria francesa, pero pensada para vivir dentro de una columna X.
\newcommand{\milcopcion}[1]{%
\makebox[5.2mm][l]{\raisebox{-.55mm}{\textcolor{milcNegro}{\milcmarca}}}%
\parbox[t]{\dimexpr\linewidth-5.2mm\relax}{\RaggedRight #1}}

% ── Recursos visuales de la guía ──────────────────────────────────────
% Declarados en el bloque `recursos:` del YAML y emitidos por el builder.
% \guiaFigura   → imagen rasterizada o PDF (foto, carta celeste, montaje).
% \guiaDiagrama → fuente TikZ/circuitikz (.tex) incrustada como vector:
%                 es la vía para circuitos y esquemáticos editables.
% #1 = ruta relativa al .tex · #2 = pie (puede ir vacío)
\newcommand{\guiaFigura}[2]{%
\begin{center}
\includegraphics[width=0.88\linewidth,height=0.38\textheight,keepaspectratio]{#1}\\[1.5mm]
{\footnotesize\itshape\textcolor{milcNegro!75}{#2}}
\end{center}
}

\newcommand{\guiaDiagrama}[2]{%
\begin{center}
\input{#1}\\[1.5mm]
{\footnotesize\itshape\textcolor{milcNegro!75}{#2}}
\end{center}
}

\begin{document}
\thispagestyle{empty}

% ── Portada ───────────────────────────────────────────────────────────
% Antes: un rectangulo de color a pagina completa. Se veia bien en pantalla,
% pero en la fotocopiadora del colegio se comia el toner y salia gris sucio.
% Ahora la banda ocupa el tercio superior y el resto va en blanco.
\begin{tikzpicture}[remember picture,overlay]
  \path[fill=milcGradePrimary]
    (current page.north west) rectangle ([yshift=-10.6cm]current page.north east);

  \node[anchor=north west,text=milcGradeText] at ([xshift=2.0cm,yshift=-1.55cm]current page.north west)
    {\sffamily\bfseries\fontsize{12}{14}\selectfont GRADO};
  \node[anchor=north west,text=milcGradeText,opacity=.85] at ([xshift=2.0cm,yshift=-2.35cm]current page.north west)
    {\sffamily\bfseries\fontsize{8.5}{10}\selectfont SERIE GUÍAS · TECNOLOGÍA E INFORMÁTICA};
  \node[anchor=north east,text=milcGradeText,opacity=.85,align=right] at ([xshift=-2.0cm,yshift=-1.55cm]current page.north east)
    {\sffamily\bfseries\fontsize{9}{12}\selectfont I.E. Sor María Juliana\\Cartago, Valle del Cauca};

  \node[anchor=west,text=milcGradeText] (gradonum) at ([xshift=1.85cm,yshift=-7.1cm]current page.north west)
    {\sffamily\bfseries\fontsize{112}{112}\selectfont 11};
  % El circulito de grado (°) solo aplica a las guias de grado («11°»); en
  % Bebras y Semillero el numero grande es un momento/modulo, no un grado.
  % Se posiciona relativo al nodo `gradonum`, no en coordenadas absolutas.
  \draw[milcGradeText,line width=1.6pt]
    ([xshift=2.0mm,yshift=-4.5mm]gradonum.north east) circle (.15cm);

  \node[anchor=west,text=milcGradeText,text width=12.3cm,align=left] at ([xshift=6.5cm,yshift=-6.6cm]current page.north west)
    {
     {\sffamily\bfseries\fontsize{9}{11}\selectfont\MakeUppercase{Período 3 · Proyecto final emprendedor}}\\[3mm]
     {\sffamily\bfseries\fontsize{21}{25}\selectfont Versión final\\del MVP ---\\la firma del aprendiz}\\[3.5mm]
     {\sffamily\bfseries\fontsize{15}{18}\selectfont Guía 29-11-TIC}
    };
\end{tikzpicture}

\vspace*{9.4cm}
\vfill

{\sffamily\bfseries\fontsize{8}{10}\selectfont\textcolor{milcGradeDark}{LO QUE TE LLEVAS HOY}}

\begin{tcolorbox}[enhanced,colback=white,colframe=milcNegro,arc=2mm,boxrule=1.6pt,
  left=5mm,right=5mm,top=4mm,bottom=4mm,before skip=3pt,after skip=12pt,
  fontupper=\RaggedRight\fontsize{11}{15.5}\selectfont]
MVP versión final integrando todo el aprendizaje del periodo (a) correcciones aplicadas a partir de los hallazgos de la sesión 5, los ajustes financieros de la sesión 7 y los compromisos éticos de la sesión 8, (b) bitácora de cambios documentando cada ajuste con su razón y referencia a la sesión que lo originó, (c) recorrido de calidad por checklist de 12 puntos verificado, (d) carta abierta de 1 página firmada con nombre y fecha que el aprendiz pueda defender públicamente ante cualquier audiencia
\end{tcolorbox}

\begin{tcolorbox}[enhanced,colback=milcGris,colframe=milcGris,arc=2mm,boxrule=0pt,
  left=5mm,right=5mm,top=3.5mm,bottom=3.5mm,after skip=12pt]
{\sffamily\bfseries\fontsize{8}{10}\selectfont\textcolor{milcNegro!55}{ESTA GUÍA ES DE}}\\[3mm]
\begin{tabularx}{\linewidth}{X p{3.2cm} p{3.2cm}}
\linead{\linewidth} & \linead{3.0cm} & \linead{3.0cm}\\[-1mm]
{\fontsize{8.5}{10}\selectfont\textcolor{milcNegro!55}{Nombre completo}} &
{\fontsize{8.5}{10}\selectfont\textcolor{milcNegro!55}{Grupo}} &
{\fontsize{8.5}{10}\selectfont\textcolor{milcNegro!55}{Fecha}}\\
\end{tabularx}
\end{tcolorbox}

\vfill

\noindent\textcolor{milcLinea}{\rule{\linewidth}{1pt}}\\[2mm]
{\sffamily\bfseries\fontsize{9.5}{12}\selectfont Autor: Dr. Álvaro Cárdenas Orozco}\\[1mm]
{\sffamily\fontsize{8.5}{11}\selectfont\textcolor{milcNegro!55}{Metodología MILC · apertura ancestral · pensamiento computacional · triángulo Dussel–estoicismo–Floridi}}

\newpage

\section*{Apertura: Versión final del MVP}

\begin{softbox}{milcGradeDark}{milcGradeSoft}{Saber ancestral}
En los talleres antiguos del Valle del Cauca y en la tradición de los oficios europeos importada al continente, había un acto ritual que marcaba el paso de aprendiz a oficial: \textbf{firmar la pieza}. El ebanista joven que llevaba tres años de aprendizaje no se llamaba todavía \emph{maestro}: era \emph{aprendiz}. El día que terminaba su primera pieza completa (una mesa, un armario, una silla), el maestro examinaba el trabajo durante una semana: revisaba uniones, lijado, ensamblaje, terminaciones. Si la pieza pasaba el examen, el maestro le permitía al aprendiz \textbf{grabar su nombre en la base con un punzón}. Esa firma era el acto público del oficio: significaba que el aprendiz se hacía responsable de la pieza para siempre, ante cualquier cliente, ante cualquier defecto futuro. Firmar implicaba 3 compromisos: (1) \textbf{defender la pieza ante cualquier crítica}; (2) \textbf{repararla si fallaba}; (3) \textbf{mantener el nombre limpio en el siguiente trabajo}. El aprendiz que rehusaba firmar no era llamado \emph{humilde}: era llamado \emph{inseguro}, y el maestro lo mandaba a rehacer la pieza hasta que pudiera firmarla con orgullo. \textbf{La firma era el verdadero examen, no el diseño.}
\end{softbox}

\begin{softbox}{milcTurquesa}{milcGris}{Saber contemporáneo}
La \textbf{versión final del MVP} es exactamente la firma del aprendiz aplicada al oficio digital. Después de 8 sesiones de trabajo (problema, validación, modelo, MVP, datos, pitch, presupuesto, ética), el aprendiz integra todos los aprendizajes en \textbf{una sola entrega defendible}. La regla profesional: \textbf{si no puedes firmar tu MVP con nombre y fecha, todavía no es versión final}. La versión final tiene 3 propiedades irrenunciables: (1) \textbf{Integra los hallazgos}: las correcciones de los datos de S5, los ajustes financieros de S7, los compromisos éticos de S8 están \emph{aplicados}, no solo \emph{prometidos}. (2) \textbf{Bitácora visible}: el aprendiz documenta cada cambio con su razón, demostrando que el MVP evolucionó por evidencia y no por capricho. (3) \textbf{Carta abierta firmada}: el aprendiz redacta una declaración pública de 1 página que pueda defender ante cualquier audiencia (familia, docentes, usuarios, asesores). La phronesis del oficio digital consiste en \textbf{firmar solo lo que se puede sostener en cualquier interrogatorio profesional}.
\end{softbox}

\begin{softbox}{milcMostaza}{milcGris}{Pregunta-puente}
\textit{¿Qué sabía el aprendiz al firmar su primera pieza ante el maestro, que el emprendedor novato olvida cuando entrega versión final sin aplicar los hallazgos? ¿Y por qué la firma es más exigente que el diseño, en cualquier oficio antiguo o contemporáneo?}
\end{softbox}

\begin{softbox}{milcVerde}{milcGris}{Saber y saber hacer}
Al terminar podrás: (1) \textbf{analizar} qué hallazgos del periodo (S5 datos, S7 finanzas, S8 ética) exigen ajustes concretos en el MVP; (2) \textbf{evaluar} cada cambio aplicado con un checklist de calidad de 12 puntos; (3) \textbf{crear} la versión final del MVP con todos los ajustes aplicados, documentada en bitácora de cambios; (4) \textbf{explicar} en una carta abierta de 1 página firmada el qué, el porqué y el para qué de tu proyecto, lista para ser leída en cualquier sustentación.
\end{softbox}



\small
\begin{tabularx}{\textwidth}{>{\bfseries}p{3.0cm}Y}
\rowcolor{milcGradePrimary!12}Autor & Dr. Álvaro Cárdenas Orozco\\
DBA & Diseño, implementación y sustentación de un proyecto de impacto local (MEN, T\&I)\\
Referentes & Dussel (analéctica · sur global) · Lean Startup · Floridi · Estoicismo\\
Producto final & MVP versión final integrando todo el aprendizaje del periodo (a) correcciones aplicadas a partir de los hallazgos de la sesión 5, los ajustes financieros de la sesión 7 y los compromisos éticos de la sesión 8, (b) bitácora de cambios documentando cada ajuste con su razón y referencia a la sesión que lo originó, (c) recorrido de calidad por checklist de 12 puntos verificado, (d) carta abierta de 1 página firmada con nombre y fecha que el aprendiz pueda defender públicamente ante cualquier audiencia\\
\end{tabularx}
\normalsize

\puente{Antes de empezar, mira el viaje completo. Llevas 8 sesiones de construcción incremental. Hoy es el penúltimo paso: \textbf{integrar todo}. La sesión 10 será sustentación pública; sin versión final lista hoy, la sustentación de mañana será improvisación. La firma del aprendiz es siempre antes del ritual público, no durante.}

\milcestacion{milcGradeDark}{Ruta MILC: así vamos a aprender}
\begin{tabularx}{\textwidth}{>{\centering\arraybackslash}X>{\centering\arraybackslash}X>{\centering\arraybackslash}X>{\centering\arraybackslash}X}
\phasecard{1}{milcVerde}{milcGris}{Escucha\\[-1mm]\small Observo, escucho y pregunto.} &
\phasecard{2}{milcTurquesa}{milcGris}{Entender\\[-1mm]\small Sistematizo y ordeno.} &
\phasecard{3}{milcMagenta}{milcGris}{Hacer\\[-1mm]\small Construyo y aplico.} &
\phasecard{4}{milcVino}{milcGris}{Cierre\\[-1mm]\small Evalúo y reflexiono.}
\end{tabularx}


\puente{Antes de teorizar, vas a hacer una \emph{lectura cruzada} de tus 3 últimas entregas: el análisis de datos (S5), el presupuesto (S7) y la autoevaluación ética (S8). Cada una declaró ajustes pendientes. Vas a listar todos los ajustes prometidos y a marcar cuáles están aplicados, cuáles a medias, cuáles pendientes.}

\milcestacion{milcVerde}{Estación 1 de 3 · Escucha}

\begin{softbox}{milcVerde}{milcGris}{Escena para escuchar}
\textbf{Actividad 1 · ANALIZA --- Inventario de ajustes pendientes por aplicar} (15 min · individual). Revisa tus 3 últimas entregas (S5 análisis de datos, S7 presupuesto, S8 ética) y lista todos los ajustes que declaraste pendientes. Para cada ajuste anota: (1) ¿está aplicado, a medias o pendiente?; (2) ¿qué se necesita para terminarlo?; (3) ¿es bloqueante para la sustentación de la próxima sesión? El objetivo es entrar a la actividad 3 con la lista cerrada de \textbf{qué falta} y \textbf{en qué orden} aplicarlo.
\end{softbox}

{\sffamily\bfseries\fontsize{8}{10}\selectfont\textcolor{milcNegro!55}{MARCA CUANDO LO HAYAS HECHO}}
\milccheckrow{Listé todos los ajustes prometidos en S5, S7, S8.}
\milccheckrow{Marqué estado (aplicado, a medias, pendiente) de cada uno.}
\milccheckrow{Identifiqué cuáles son bloqueantes para la sustentación.}

\begin{infoband}{milcVerde}


\textbf{Lo que va al cuaderno (Actividad 1 · ANALIZA).} Título: «Actividad 1 --- Ajustes pendientes». \textbf{Formato:} tabla 3+ filas × 4 columnas (Ajuste~|~Sesión origen~|~Estado~|~Bloqueante sí/no). \textbf{Sabes que terminaste cuando} tienes claro qué orden seguirás para aplicar los ajustes esta semana.
\end{infoband}

\puente{Ya tienes la lista de ajustes. Ahora vas a entender la \textbf{anatomía de la versión final}: las 3 piezas obligatorias (MVP ajustado, bitácora de cambios, carta abierta firmada), el \textbf{checklist de 12 puntos} de calidad, y los 6 errores que delatan un MVP \emph{``casi final''}.}

\milcestacion{milcTurquesa}{Estación 2 de 3 · Entender}

\begin{softbox}{milcTurquesa}{milcGris}{Lo que vas a entender}
Una versión final sin método de cierre es un MVP cansado, no terminado. Para que tu firma de hoy tenga peso necesitas la \textbf{anatomía profesional del cierre}: 3 piezas, 12 puntos de calidad, 6 errores típicos del aprendiz que firma sin haber pulido.
\end{softbox}

\milcpaso{1}{milcTurquesa}{La versión final se descompone en 3 piezas: (1) \textbf{MVP ajustado}: la versión actualizada con los ajustes aplicados (no solo prometidos); cambios visibles en captura, valor o respuesta. (2) \textbf{Bitácora de cambios}: documento corto donde cada ajuste tiene una fila (fecha, qué cambió, por qué, referencia a la sesión que lo originó). (3) \textbf{Carta abierta firmada}: declaración pública de 1 página estructurada en 3 párrafos (qué hace el proyecto, por qué importa, qué se compromete a entregar el aprendiz) cerrada con nombre y fecha.}
\milcpaso{2}{milcTurquesa}{El \textbf{checklist de 12 puntos} de calidad cubre 4 dimensiones, 3 puntos cada una: (a) \textbf{Producto}: (1) el MVP funciona end-to-end sin fricciones graves, (2) el copy de cada pantalla es claro y honesto, (3) los enlaces internos no están rotos. (b) \textbf{Datos}: (4) la captura registra correctamente, (5) la base de datos está organizada y consultable, (6) los datos pueden eliminarse a petición del usuario. (c) \textbf{Finanzas}: (7) los precios están visibles si aplica, (8) los costos del MVP están actualizados, (9) el punto de equilibrio mes 3 es alcanzable según escenario realista. (d) \textbf{Ética}: (10) Política de Privacidad accesible desde el MVP, (11) Términos de Uso accesibles, (12) ajuste ético declarado en S8 está aplicado.}
\milcpaso{3}{milcTurquesa}{Lo esencial cabe en una pregunta: \emph{¿firmaría este MVP con mi nombre y mi fecha hoy mismo?}. Si la respuesta es no, todavía no es versión final. La phronesis del aprendiz consiste en \textbf{no firmar lo que no se puede defender}. Firmar antes de tiempo erosiona el oficio; no firmar nunca por inseguridad también.}
\milcpaso{4}{milcTurquesa}{Algoritmo: (1) aplica los ajustes pendientes en orden de bloqueo; (2) recorre el checklist de 12 puntos marcando OK o pendiente; (3) corrige los pendientes hasta tener 12/12; (4) llena la bitácora de cambios con cada ajuste; (5) redacta la carta abierta de 1 página; (6) firma con nombre y fecha; (7) léela en voz alta para verificar que la podrías defender ante cualquier audiencia.}

\begin{drawbox}{milcTurquesa}{Anatomía de la versión final}
\textbf{MVP ajustado:} con cambios aplicados, no solo prometidos. \\[1mm] \textbf{Bitácora de cambios:} fecha + cambio + razón + sesión origen. \\[1mm] \textbf{Checklist de calidad:} 12 puntos en 4 dimensiones. \\[1mm] \textbf{Carta abierta firmada:} 1 página + nombre + fecha. \\[1mm] \textbf{Criterio interno:} ¿firmaría con orgullo hoy?
\end{drawbox}

\begin{softbox}{milcMostaza}{milcGris}{Errores comunes y su corrección}
(1) \textbf{Prometer ajustes y no aplicarlos}: la versión final sin ajustes aplicados es la misma de la sesión 4. (2) \textbf{Bitácora vacía o decorativa}: cada fila debe tener razón y referencia, no descripción genérica. (3) \textbf{Checklist marcado sin revisar}: marcar OK sin probar es autoengaño. (4) \textbf{Carta abierta genérica}: \emph{``este proyecto busca mejorar la vida de las personas''} sin contenido específico es relleno. (5) \textbf{Firma sin convicción}: firmar dudando es señal de trabajo incompleto. (6) \textbf{Carta sin fecha}: la fecha es parte de la firma profesional, no detalle.
\end{softbox}

\begin{infoband}{milcTurquesa}


\textbf{Lo que va al cuaderno (Actividad 2 · EVALÚA).} Título: «Actividad 2 --- Anatomía del cierre». \textbf{Formato:} (a) las 3 piezas anatómicas explicadas; (b) checklist de 12 puntos listo para marcar; (c) los 6 errores con marca del que más temes. \textbf{Sabes que terminaste cuando} entiendes que la firma es el examen real, no el diseño.
\end{infoband}

\puente{Tienes el método. Toca aplicarlo: aplicar los ajustes pendientes, llenar la bitácora con cada cambio y su razón, recorrer el checklist de 12 puntos, redactar la carta abierta y firmarla con nombre y fecha.}

\milcestacion{milcMagenta}{Estación 3 de 3 · Hacer}

\begin{softbox}{milcMagenta}{milcGris}{Cómo vas a construir tu producto}
\textbf{Actividad 3 · CREA + EXPLICA --- Versión final firmada} (45 min en sesión, 3-4 horas de trabajo distribuido en la semana · individual). Hora de aplicar y firmar. Aplicas los ajustes pendientes, recorres el checklist hasta 12/12, llenas la bitácora, redactas la carta y firmas.
\end{softbox}

\milcpaso{1}{milcMagenta}{Tu versión final se produce en 6 pasos: (1) aplicar los ajustes pendientes en orden de bloqueo (los bloqueantes primero); (2) recorrer el checklist marcando cada punto con OK o pendiente; (3) corregir los pendientes hasta 12/12; (4) escribir la bitácora con 1 fila por cambio aplicado; (5) redactar la carta abierta de 1 página en 3 párrafos (qué, por qué, compromiso); (6) firmar con nombre y fecha.}
\milcpaso{2}{milcMagenta}{Sigue el patrón \textbf{``aplicación sobre declaración''}: en la versión final, no vale decir \emph{``haré X''}; debe decir \emph{``hice X el día Y''}. La diferencia entre versión final y borrador es esa: la versión final es performativa (lo que dice, ya ocurrió), el borrador es declarativa (lo que dice, todavía no ocurre).}
\milcpaso{3}{milcMagenta}{La carta abierta tiene estructura mínima de 3 párrafos: (a) \textbf{Qué}: \emph{``Este proyecto se llama X, atiende el problema de Y, está dirigido al segmento Z, y se sostiene con el modelo W''}. (b) \textbf{Por qué}: \emph{``Lo construí porque [relación personal con el problema]. Aprendí en 10 sesiones que [hallazgo más importante]. Las personas que validaron el proyecto fueron [referencia a las 5 entrevistas]''}. (c) \textbf{Compromiso}: \emph{``Me comprometo a [cumplir lo prometido en Política, Términos, Autoevaluación]. Si el proyecto crece, mantendré [los compromisos éticos]. Si falla, [cómo asumo responsabilidad]''}. Firma con nombre completo + fecha.}
\milcpaso{4}{milcMagenta}{Algoritmo final: (1) aplica ajustes; (2) checklist; (3) corrige; (4) bitácora; (5) carta; (6) firma; (7) lee la carta en voz alta y verifica si la podrías defender ante alguien que no te conoce; (8) si dudas, vuelve a 1.}

\begin{drawbox}{milcMagenta}{Checklist final de los 12 puntos}
\checkbox MVP end-to-end sin fricciones graves. \\[1mm] \checkbox Copy claro y honesto en cada pantalla. \\[1mm] \checkbox Enlaces internos funcionando. \\[1mm] \checkbox Captura de datos registrando bien. \\[1mm] \checkbox Base de datos organizada y consultable. \\[1mm] \checkbox Datos eliminables a petición del usuario. \\[1mm] \checkbox Precios visibles (si aplica). \\[1mm] \checkbox Costos actualizados. \\[1mm] \checkbox Punto de equilibrio alcanzable según escenario realista. \\[1mm] \checkbox Política de Privacidad accesible desde el MVP. \\[1mm] \checkbox Términos de Uso accesibles. \\[1mm] \checkbox Ajuste ético declarado en S8 aplicado.
\end{drawbox}

\begin{softbox}{milcMagenta}{milcGris}{Plantilla de trabajo}
\textbf{Bitácora.} \textit{Fecha + cambio + razón + sesión origen.} \\[1mm] \textbf{Checklist.} \textit{12 puntos marcados con OK.} \\[1mm] \textbf{Carta párrafo 1 (qué).} \textit{Nombre del proyecto + problema + segmento + modelo.} \\[1mm] \textbf{Carta párrafo 2 (por qué).} \textit{Relación personal + hallazgo principal + validación.} \\[1mm] \textbf{Carta párrafo 3 (compromiso).} \textit{Compromisos éticos + responsabilidad si falla.} \\[1mm] \textbf{Firma.} \textit{Nombre completo + fecha.}
\end{softbox}

\begin{infoband}{milcMagenta}


\textbf{Lo que va al cuaderno (Actividad 3 · CREA + EXPLICA).} Título: «Actividad 3 --- Versión final firmada». \textbf{Formato:} (a) bitácora completa de cambios; (b) checklist de 12 puntos verificado; (c) carta abierta de 1 página firmada con nombre y fecha. \textbf{Extensión:} 2-3 páginas. \textbf{Sabes que terminaste cuando} podrías mostrar tu MVP final a cualquier persona y leer la carta sin dudar.
\end{infoband}

\puente{Lo que sigue resume \emph{qué} entregas hoy: MVP versión final + bitácora de cambios + checklist verificado + carta abierta firmada. El criterio principal: que el aprendiz pueda decir con sinceridad \emph{``firmo este proyecto y lo defiendo''} ante cualquier audiencia. Si hay duda al firmar, falta trabajo, no falta valentía.}

\milcestacion{milcMostaza}{Producto de la sesión}
\textbf{MVP versión final + bitácora + checklist 12/12 + carta abierta firmada}.
\begin{softbox}{milcMostaza}{milcGris}{Criterios de calidad}
(1) Ajustes pendientes de S5, S7, S8 aplicados. (2) Bitácora con 1 fila por cambio (fecha, qué, por qué, sesión origen). (3) Checklist de 12 puntos completado en las 4 dimensiones. (4) Carta abierta de 1 página en 3 párrafos. (5) Firma con nombre completo y fecha. (6) Capacidad de defender el MVP ante audiencia desconocida.
\end{softbox}

\puente{Antes de cerrar, mira la versión final desde las cinco dimensiones humanas. Firmar el oficio terminado es respeto por el proceso entero; entregar sin aplicar hallazgos es traicionar las 8 sesiones anteriores.}

\milcestacion{milcVino}{Evaluación liberadora · cinco dimensiones}

\begin{softbox}{milcVino}{milcGris}{Sentido de la evaluación}
Evaluar no es solo revisar una nota. Es reconocer qué aprendí, qué cambió en mí, qué emociones manejé, cómo me relaciono con la ciudad y la comunidad, y qué le dejaré a quien viene detrás.
\end{softbox}

\textbf{Mi reflexión liberadora en cinco dimensiones.} Completa cada frase iniciada:

\small
\begin{tabularx}{\textwidth}{>{\bfseries\RaggedRight\arraybackslash}p{3.4cm}Y>{\RaggedRight\arraybackslash}p{4.6cm}}
\rowcolor{milcVino}\color{white}Dimensión & \color{white}\bfseries Mi respuesta (frame starter) & \color{white}\bfseries Algo que descubrí\\
1. Personal & Lo que más cambió en mí fue \linead{6.5cm} & \linead{4.2cm}\\
2. Emocional & La emoción más fuerte fue \linead{2.5cm} y la regulé \linead{3cm} & \linead{4.2cm}\\
3. Ciudadana & En democracia esto importa porque \linead{6cm} & \linead{4.2cm}\\
4. Local (Cartago) & En mi barrio sería útil para \linead{6.5cm} & \linead{4.2cm}\\
5. Intergeneracional & A mi abuela/abuelo le diría \linead{6.5cm} & \linead{4.2cm}\\
\end{tabularx}
\normalsize

\begin{infoband}{milcVino}
\textbf{Para la próxima sustentación.} Vuelve a esta tabla en seis meses. Verás que las respuestas cambian — no porque lo aprendido cambie, sino porque tú cambias. La evaluación liberadora es una conversación contigo a través del tiempo.
\end{infoband}


\puente{Y para terminar, tres voces sobre la firma y la responsabilidad. Dussel sobre los productos que asumen al usuario, Marco Aurelio sobre la coherencia del oficio bien hecho, Floridi sobre la responsabilidad digital como ética del autor.}

\milcestacion{milcGradeDark}{Triángulo de pensamiento · cierre filosófico}

\begin{softbox}{milcVino}{milcGris}{Dussel · lente del nosotros}
\textit{Firmar un proyecto es asumir al usuario como prójimo, no como cliente abstracto.}

Cuando firmas con nombre y fecha, declaras públicamente que tu proyecto tiene una persona responsable detrás, no una entidad anónima. La filosofía de la liberación pide ese acto: nombrarse para asumir responsabilidad por el otro. Las empresas grandes a menudo se esconden detrás de razones sociales y términos legales para diluir responsabilidad; el aprendiz que firma su MVP con su nombre propio hace lo contrario. La phronesis del oficio consiste en \textbf{no esconderse detrás del sistema cuando se firma}.

\textbf{¿Estoy dispuesto a poner mi nombre real y mi fecha real en la carta abierta? ¿O quisiera firmar bajo seudónimo para no asumir la responsabilidad?}
\end{softbox}
\linead{\linewidth}\\[1mm]
\linead{\linewidth}

\begin{softbox}{milcMostaza}{milcGris}{Estoicismo · Marco Aurelio · cuidado interior}
\textit{El oficio bien hecho es la única firma que dura; los atajos firman con tinta que se borra.}

Es tentador firmar la versión final sin haber aplicado los ajustes, esperando que la audiencia no se dé cuenta. La disciplina estoica recomienda lo opuesto: \emph{firma solo lo que se puede sostener bajo examen}. La coherencia entre lo prometido y lo entregado es virtud profesional irrenunciable. Quien firma con descuido pierde la única autoridad que la vida le ofrece: el respeto del oficio bien terminado.

\textbf{Si alguien revisara mi MVP punto por punto contra mis 3 últimas entregas, ¿encontraría coherencia entre lo prometido y lo hecho?}
\end{softbox}
\linead{\linewidth}\\[1mm]
\linead{\linewidth}

\begin{softbox}{milcTurquesa}{milcGris}{Floridi · lente de la infoesfera}
\textit{Firmar digitalmente con responsabilidad es la nueva forma de la autoría ética en la infosfera.}

En la era de la información, los proyectos digitales suelen ser anónimos o colectivos diluidos; identificar al autor es ya un acto de ética. La carta abierta firmada te entrena en una práctica que la cultura digital contemporánea ha erosionado: la autoría asumida. Quien firma su MVP escolar con su nombre y fecha hoy se entrena para firmar sus decisiones profesionales mañana. La autoría es la primera línea de la responsabilidad digital.

\textbf{¿Mi carta firmada me identifica como autor responsable o me diluye en lenguaje impersonal que esconde quién decide qué?}
\end{softbox}
\linead{\linewidth}\\[1mm]
\linead{\linewidth}

\begin{infoband}{milcNegro}
\textbf{Nota para el docente.} Las tres formulaciones del triángulo son la síntesis del autor de la guía, no citas textuales de los pensadores. Si vas a citarlos en clase, remite a la obra original.
\end{infoband}

\milcestacion{milcVino}{Mi autoevaluación · marca una casilla por fila}

{\small
\setlength{\extrarowheight}{2.2pt}
\begin{tabularx}{\textwidth}{>{\RaggedRight\arraybackslash\bfseries}p{3.0cm}YYY}
\rowcolor{milcVino}\color{white}\bfseries Criterio & \color{white}\bfseries Lo logré & \color{white}\bfseries En proceso & \color{white}\bfseries Necesito apoyo\\[.6mm]
Comprensión del tema & \milcopcion{Explico las ideas con mis palabras.} & \milcopcion{Reconozco las ideas, falta precisión.} & \milcopcion{Necesito repasar lo básico.}\\[.8mm]
\rowcolor{milcGris}Pensamiento computacional & \milcopcion{Usé los pilares al armar mi producto.} & \milcopcion{Usé algunos pilares.} & \milcopcion{Necesito releer la estación 2.}\\[.8mm]
Aplicación y producto & \milcopcion{Mi producto cumple los criterios.} & \milcopcion{Está armado, con ajustes pendientes.} & \milcopcion{Aún está en construcción.}\\[.8mm]
\rowcolor{milcGris}Reflexión 5 dimensiones & \milcopcion{Respondí cada una con honestidad.} & \milcopcion{Respondí algunas con detalle.} & \milcopcion{Mis respuestas son muy generales.}\\[.8mm]
Triángulo de pensamiento & \milcopcion{Conecté las 3 voces con mi trabajo.} & \milcopcion{Conecté al menos una.} & \milcopcion{No las relaciono con mi proceso.}\\[.8mm]
\end{tabularx}}

\vspace{3mm}
\textbf{Hoy entendí que:}\\[1mm]
\linead{\linewidth}\\[1mm]
\linead{\linewidth}

\textbf{Mi compromiso para la próxima sesión es:}\\[1mm]
\linead{\linewidth}\\[1mm]
\linead{\linewidth}

\begin{softbox}{milcMostaza}{milcGris}{Cita de despedida · Marco Aurelio}
\textit{``El obstáculo en el camino se vuelve el camino. Lo que estorba al camino se vuelve camino.''}\\[1mm]
Cada error de hoy, bien observado, se convierte en pieza del camino que pisarás mañana.
\end{softbox}

\small
\begin{tabularx}{\textwidth}{>{\bfseries}p{3.6cm}Y}
\rowcolor{milcGradeDark!12}Nombre completo: & \linead{10cm}\\
Fecha: & \linead{4cm} \quad Curso/grupo: \linead{4cm}\\
Mi firma: & \linead{9cm}\\
\end{tabularx}
\normalsize

\begin{infoband}{milcGradeDark}
\textbf{Para la próxima sesión.} Llega con el MVP versión final, bitácora, checklist verificado y carta abierta firmada. En la sesión 10 vas a hacer la \textbf{sustentación pública}: 15 minutos frente a la audiencia (compañeros, docentes, posiblemente padres y asesores externos) + preguntas + manifiesto final del grado 11. Sin versión final firmada hoy, la sustentación de mañana será defensa de algo todavía en construcción.
\end{infoband}

\end{document}
